\section{Funzionamento dell'applicazione}
\label{sec:funzionamento}

\subsection{Schermata principale}

L'applicazione all'avvio mostrerà la \textbf{Home Screen}, che fornisce un collegamento diretto a tutte le funzionalità offerte dall'app.

La schermata è suddivisa in tre parti principali: la prima mostra il \textbf{Device ID} del dispositivo, al lato l'icona in verde conferma la corretta autenticazione del dispositivo; 

poco sotto ci sono i bottoni di navigazione per le schermate dell'app: il primo in verde per iniziare un trasfermento porterà alla schermata dove inizierà la discovery dei dispositivi circostanti,

i due tasti sotto portano invece portano rispettivamente alla history delle transazioni e alle opzioni, dove l'utente ha la possibilità di abilitare il synch con il backend e gestire i gruppi.

nella parte sottostante sono mostrate le ultime tre transazioni effettuate, ciascuna mostrata in una card con il nome del file, la dimensione, la direzione (invio o ricezione) e il timestamp.

\vspace{12pt}
\begin{figure}[h]
    \centering
    \includegraphics[width=0.42\textwidth]{home_page.png}
    \caption{Home Screen con le ultime transazioni recenti}
    \label{fig:home}
\end{figure}

\subsection{Navigazione}

\subsubsection*{Discovery e invio file}

Premendo il tasto verde nella Home l'utente accede alla \textbf{Discovery Screen}. L'app avvia automaticamente la scansione Wi-Fi Direct e mostra i dispositivi rilevati nelle vicinanze. Per avviare un trasferimento l'utente seleziona prima il dispositivo destinatario e poi sceglie il file dal file picker di sistema; a quel punto inizierà il traferimento, facendo comparire la notifica persistente di avanzamento.

\vspace{12pt}
\begin{figure}[h]
    \centering
    \includegraphics[width=0.65\textwidth]{pairing_page.png}
    \caption{Discovery Screen con i dispositivi rilevati nelle vicinanze}
    \label{fig:discovery}
\end{figure}

\pagebreak
\subsubsection*{History}

La \textbf{History Screen} elenca tutte le transazioni registrate, sia locali che sincronizzate dal gruppo. Ogni \texttt{HistoryCard} è collassabile: inizialmente mostra le info principali come nome del file e il timestamp; espandendola compaiono anche il Device ID del peer e, se disponibile, la posizione in cui è avvenuto il trasferimento.

\vspace{12pt}
\begin{figure}[h]
    \centering
    \includegraphics[width=0.42\textwidth]{history_page.png}
    \caption{History Screen con le transazioni registrate}
    \label{fig:history}
\end{figure}

\pagebreak
\subsubsection*{Impostazioni e gruppi}

La \textbf{Settings Screen} permette di abilitare la sincronizzazione con il backend e di gestire i gruppi. Attivandolo, l'app registra il dispositivo sul backend ottenendo il token. Nella sezione gruppi l'utente può creare un nuovo gruppo oppure unirsi a uno esistente inserendo il codice; al join, tutte le entry locali vengono caricate sul backend con il codice del gruppo, rendendole visibili agli altri membri.

\vspace{12pt}
\begin{figure}[h]
    \centering
    \begin{minipage}[t]{0.42\textwidth}
        \centering
        \includegraphics[width=\textwidth]{settings_page.png}
        \caption{Settings Screen}
        \label{fig:settings}
    \end{minipage}
    \hfill
    \begin{minipage}[t]{0.42\textwidth}
        \centering
        \includegraphics[width=\textwidth]{settings_page_group_created.png}
        \caption{Gruppo creato con codice condivisibile}
        \label{fig:settings-group}
    \end{minipage}
\end{figure}

\clearpage
